\section{Experiments}
\label{sec:experiments}

\subsection{Experimental Setup}
\label{sec:setup}

\paragraph{Dataset.} We evaluate on the GlobalMMLU test split \citep{singh2024globalmmlu}
across 20 typologically diverse languages grouped into three resource tiers: \textsc{high} (zh, de, es, fr, ja, ko, pt, ru; 8 languages), \textsc{med} (ar, id, nl, pl, tr, uk, vi; 7 languages), and \textsc{low} (bn, hi, sw, te; 4 languages). The test split is partitioned 70/30 via stratified sampling into \texttt{probe\_train} and \texttt{eval} portions (\S\ref{sec:framework}). English correctness labels from \texttt{probe\_train} serve as the sole supervision; all reported metrics use \texttt{eval} exclusively.

\paragraph{Models.} We evaluate two architecture families at two scales. Primary models (P0): Qwen3-235B-A22B (GPTQ-Int4, 94 layers, ${\sim}$125\,GB VRAM under 4-GPU tensor parallelism) and Llama~4 Scout 17B-16E (W4A16 weight quantization, 48 layers, mixture-of-experts with 16 experts). Pilot models (MVP): Qwen2.5-7B-Instruct (INT8, 28 layers) and Llama-3.1-8B-Instruct (INT8, 32 layers). The pilot pair spans both architecture families and validates methodology and hyperparameter choices before committing P0 GPU budget.

\paragraph{Activation Extraction.} We register forward hooks on the post-layer-norm residual stream at every 4th layer plus the final layer, yielding 8 probed layers for the 28-layer Qwen2.5-7B, 9 for the 32-layer Llama-3.1-8B, 13 for the 48-layer Llama~4 Scout, and 25 for the 94-layer Qwen3-235B.

\paragraph{Probing Protocol.} Logistic regression classifiers (liblinear solver, L2 regularization, $C{=}1.0$, \texttt{max\_iter}{=}2000) are trained on English \texttt{probe\_train} correctness labels. Standard probes are trained on full-dimensional activations (dimensions vary by architecture, e.g., 3584-d for 7B-scale models). Main probes are evaluated on the held-out \texttt{eval} split (30\%), with the same item partition shared across all languages to preserve parallelism. For control task analysis (\S\ref{sec:controls}), activations are reduced to 100 dimensions via PCA and evaluated with 5-fold stratified cross-validation. We compare five transfer conditions as described in \S\ref{sec:transfer}: zero-shot, mean-shift, Procrustes, few-shot with 25 and 100 target-language labeled examples, and CaCLR.

\paragraph{Evaluation Metrics.} AUROC serves as the primary discrimination metric. We report Soft-Binned ECE (SB-ECE) \citep{karandikar2021sbece} for calibration quality and Spearman~$\rho$ for the depth-transfer correlation analysis (\S\ref{sec:geometric}). Standard ECE and Brier score are reported in the Appendix.

\paragraph{Statistical Testing.} All reported metrics include bootstrap confidence intervals (1000 resamples). Method comparisons use paired permutation tests; we report Cohen's~$d$ effect sizes for key comparisons. The CKA-transfer analysis (\S\ref{sec:geometric}) employs partial correlation controlling for layer depth. All $p$-values are two-sided with significance threshold $\alpha{=}0.05$.

\paragraph{Hardware.} All experiments run on 4$\times$ NVIDIA RTX PRO 6000 Blackwell Server Edition (96\,GB each). The Qwen3-235B model requires all four GPUs under tensor parallelism; pilot and Scout models each fit on a single device. % [TODO: total GPU-hours estimate]

\subsection{Main Transfer Results}
\label{sec:main_results}

\subsection{Pareto Frontier Analysis}
\label{sec:pareto_results}

\subsection{CKA Geometric Analysis}
\label{sec:cka_results}

\subsection{Decomposition Analysis}
\label{sec:decomp_results}

\subsection{Empirical Findings}
\label{sec:empirical}

\subsection{CaCLR Analysis}
\label{sec:caclr_results}
